\documentclass[twocolumn]{aastex631}
\usepackage{amsmath}
\usepackage{booktabs}
\shorttitle{SDSS denominator/proxy atlas}
\shortauthors{NebulaMind}
\begin{document}

\title{Supplementary SDSS Denominator and Proxy Atlas for Galaxy-Evolution Follow-up}
\author{NebulaMind Research Autopilot}
\affiliation{Public SDSS DR17 data only}

\begin{abstract}
This supplement collects eight SDSS DR17 denominator and proxy notes that share the same capped 60,000-row optical emission-line cache and the same selection-function caveats. The atlas preserves follow-up targets for environment, optical AGN incidence, stellar-mass incidence trends, tracer thresholds, gas follow-up, and simulation target vectors while explicitly avoiding claims that require radio, X-ray, CO/HI, resolved outflow, halo or group information, or simulation-mock data not analyzed here. It is a single follow-up atlas, not eight independent causal-feedback papers.
\end{abstract}

\keywords{galaxies: evolution --- surveys --- catalogs --- methods: observational --- methods: statistical}

\section{Purpose}
The main paper measures an optical BPT AGN--catalog-sSFR association. These eight topics are distinct: each is scientifically useful as a denominator or target-vector definition, but each lacks at least one core physical observable required by its original proposal. The BPT language and catalog-backbone language here follow the same SDSS/MPA-JHU-style value-added tables and standard demarcations as the flagship \citep{sdssdr17,brinchmann2004,york2000,baldwin1981,kewley2001,kauffmann2003bpt,kewley2006,stasinska2008,stasinska2015}. Keeping them in one supplement prevents overclaiming and gives future work a single checklist of what still must be added.

\section{Shared denominator}
The atlas uses the same cached public-data backbone as the main paper: 60,000 cached rows from a strict public four-line S/N$\geq3$ parent of 249,917 rows, i.e. 24.0\% cached coverage. The four-line selection is sSFR-dependent and the cache is capped and non-random, so all counts and fractions are conditional denominators rather than population-complete measurements. The row-level stellar masses and catalog sSFR values are taken from the public MPA-JHU-style \texttt{galSpecExtra} table after the same SDSS joins used in the flagship \citep{sdssdr17,brinchmann2004,york2000}.

The eight subsections below are intentionally parallel: each one states the observed optical denominator or target vector, then lists the missing observables that a future multiwavelength or simulation-based test would have to add before any physical inference can be made.

\begin{deluxetable*}{lrrr}
\tabletypesize{\scriptsize}
\tablecaption{Selection cascade shared by the atlas.\label{tab:supp-selection}}
\tablehead{\colhead{Selection stage} & \colhead{Public DR17 rows} & \colhead{Cached rows} & \colhead{Retention vs. spectro-z parent}}
\startdata
SpecObj GALAXY, 0.02<z<0.12 & 501,060 & -- & 1.000 \\
plus galSpecInfo/PhotoObj/galSpecExtra and mass/sSFR bounds & 416,554 & -- & 0.831 \\
plus galSpecLine join & 416,554 & -- & 0.831 \\
four BPT lines positive with positive errors & 373,445 & 60,000 & 0.745 \\
four BPT lines S/N>=3 & 249,917 & 60,000 & 0.499 \\
four BPT lines S/N>=5 & 176,523 & 42,446 & 0.352 \\
four BPT lines S/N>=10 & 91,768 & 22,311 & 0.183 \\
\enddata
\end{deluxetable*}

\section{Atlas notes}

\subsection{SDSS density proxy for low-sSFR incidence}
The follow-up goal here is to isolate an environmental denominator that can later be joined to group catalogs and halo masses. The nearest-neighbour density proxy adds low-sSFR incidence information beyond stellar mass in the SDSS emission-line sample. The SDSS emission-line denominator contains 60,000 galaxies with an internally computed 10th-neighbour density proxy. The high-density quartile has a low-sSFR emission-line fraction of 0.230 (3,456/15,000), while the low-density quartile has 0.181 (2,710/15,000). The bootstrap high-minus-low interval is [0.041, 0.059], and a linear probability model adjusted for log stellar mass and redshift gives a high-density coefficient of 0.032 +/- 0.004. This is a denominator-level environmental diagnostic; group catalogues, robust central/satellite labels, halo masses, morphology, and multi-redshift selection functions are still needed for a physical environmental interpretation \citep{peng2010,wetzel2013,dekel2006}.

\begin{figure}
\centering
\includegraphics[width=\columnwidth]{../figures/topic-01.pdf}
\caption{SDSS optical denominator/proxy diagnostic for m1\_rp2\_environment\_quenching. This is a follow-up target definition or baseline, not a physical-feedback measurement.}
\label{fig:m1-rp2-environment-quenching}
\end{figure}


\subsection{Optical-AGN denominator for maintenance-heating follow-up}
The follow-up goal here is to identify the optical-AGN duty-cycle denominator that radio and X-ray data would need to test maintenance heating. Among massive, low-sSFR SDSS emission-line galaxies, the optical AGN fraction can be used as a denominator for X-ray and radio maintenance-heating follow-up. The massive subset (\(\log M_\star \geq 10.8\)) contains 9,298 emission-line galaxies, of which 5,695 are low-sSFR by the pilot threshold. The optical BPT AGN fraction is 0.430 in the massive subset and 0.607 among massive low-sSFR objects. This provides an optical duty-cycle denominator for X-ray and radio follow-up, not a heating-to-cooling measurement. The missing observables are X-ray cavity or cooling-luminosity measurements, radio jet powers, halo-selected parent catalogues, and nondetection modelling \citep{best2005,heckmanbest2014,fabian2012,mcnamara2007,lamassa2013}.

\begin{figure}
\centering
\includegraphics[width=\columnwidth]{../figures/topic-02.pdf}
\caption{SDSS optical denominator/proxy diagnostic for m1\_rp3\_maintenance\_heating. This is a follow-up target definition or baseline, not a physical-feedback measurement.}
\label{fig:m1-rp3-maintenance-heating}
\end{figure}


\subsection{SDSS high-excitation AGN denominator for outflow tests}
The follow-up goal here is to isolate the optical-AGN denominator that resolved kinematics would need to test escape versus recycling. The SDSS high-excitation optical-AGN denominator identifies how many systems would need resolved kinematics to test escape versus recycling. High-excitation optical AGN candidates number 4,440 of 60,000 emission-line galaxies (0.074). Their median \(\log {\rm sSFR}\) is -11.53, compared with -10.14 for the full denominator. SDSS does not measure escape velocity or multiphase outflow velocities here; the note supplies a denominator for resolved follow-up rather than an escape or recycling result. The missing observables are resolved outflow velocities, halo potentials, molecular, ionized, and neutral gas phases, and CGM recycling tracers \citep{veilleux2005,cicone2014,carniani2017,fiore2017,lamassa2013}.

\begin{figure}
\centering
\includegraphics[width=\columnwidth]{../figures/topic-03.pdf}
\caption{SDSS optical denominator/proxy diagnostic for m2\_p1\_outflow\_escape\_recycling. This is a follow-up target definition or baseline, not a physical-feedback measurement.}
\label{fig:m2-p1-outflow-escape-recycling}
\end{figure}


\subsection{Environment proxy for optical AGN in massive SDSS hosts}
The follow-up goal here is to define the environment-stratified optical denominator that future radio and X-ray work could test. The local-density proxy is correlated with the optical AGN fraction in massive SDSS hosts and motivates environment-stratified radio and X-ray follow-up. Among massive hosts, the high-density quartile has an optical AGN fraction of 0.509, while the low-density quartile has 0.367. The bootstrap high-minus-low interval is [0.112, 0.170]. This is an optical/environment denominator for radio-jet coupling work; it does not measure radio jet power or coupling efficiency. The missing observables are radio jet morphology and age, cavity or shock energetics, hot-gas density, and calibrated jet-power estimates \citep{best2005,mcnamara2007,heckmanbest2014}.

\begin{figure}
\centering
\includegraphics[width=\columnwidth]{../figures/topic-04.pdf}
\caption{SDSS optical denominator/proxy diagnostic for m2\_p2\_radio\_jet\_environment. This is a follow-up target definition or baseline, not a physical-feedback measurement.}
\label{fig:m2-p2-radio-jet-environment}
\end{figure}


\subsection{Stellar-mass distribution of low-sSFR and optical AGN incidence}
The follow-up goal here is to pin down the mass bin where a future gas-inclusive study should look for an incidence change. At what stellar-mass scale do the low-sSFR emission-line fraction and optical AGN incidence rise in the same SDSS denominator? The first stellar-mass bin with low-sSFR fraction above 0.5 is \(\log(M_\star/M_\odot) \in [11.0,12.5]\). The optical AGN fraction peaks in the 11.0-12.5 bin at 0.520. This is an optical distribution diagnostic; gas fractions and baryon deficits are needed before assigning any physical meaning to the apparent incidence change. The missing observables are gas fractions, baryon deficits, halo masses, stellar-feedback observables, and high-redshift extensions.
The same binning is therefore best treated as a population-distribution diagnostic, not a statement about a physical transition mass for individual galaxies \citep{peng2010,wetzel2013,dekel2006}.

\begin{figure}
\centering
\includegraphics[width=\columnwidth]{../figures/topic-05.pdf}
\caption{SDSS optical denominator/proxy diagnostic for m2\_p3\_feedback\_transition\_mass. This is a follow-up target definition or baseline, not a physical-feedback measurement.}
\label{fig:m2-p3-feedback-transition-mass}
\end{figure}


\subsection{Common-denominator optical tracer census in SDSS}
The follow-up goal here is to compare optical tracer choices against one shared denominator before any multiphase census is attempted. How strongly do simple optical tracer definitions change the inferred AGN or feedback-candidate prevalence in one common SDSS denominator? Within the same 60,000-galaxy denominator, simple optical tracer definitions produce prevalence from 0.136 to 0.418. The widest-to-narrowest prevalence ratio is 3.1 before adding molecular, neutral, or X-ray or radio phases. This demonstrates why a common-denominator multiphase census is required; it does not measure molecular or neutral outflow rates. The missing observables are ionized, molecular, neutral, and X-ray or radio tracers measured over the same parent denominator and aperture model \citep{xcoldgass2017,xgass2018,cicone2014,carniani2017,fiore2017,veilleux2005}.

\begin{figure}
\centering
\includegraphics[width=\columnwidth]{../figures/topic-06.pdf}
\caption{SDSS optical denominator/proxy diagnostic for m3\_p1\_multiphase\_census. This is a follow-up target definition or baseline, not a physical-feedback measurement.}
\label{fig:m3-p1-multiphase-census}
\end{figure}


\subsection{Optical denominator for gas-fraction versus efficiency tests}
The follow-up goal here is to define the denominator for CO/HI gas-fraction and depletion-time follow-up. How many massive low-sSFR or transitioning SDSS galaxies with valid emission-line measurements are available as a denominator for CO gas-fraction and depletion-time follow-up? The massive low-sSFR denominator contains 6,729 galaxies in the SDSS emission-line sample. Its optical BPT AGN fraction is 0.549, and the median H-alpha luminosity proxy is 40.06. The median H-alpha luminosity proxy is 0.66 dex lower than in massive star-forming emission-line galaxies. SDSS optical data cannot distinguish molecular-gas depletion from reduced star-formation efficiency; this note identifies the CO/HI follow-up denominator and optical baseline. The missing observables are CO or dust-based molecular gas masses, aperture-matched SFRs, morphology, and environment labels \citep{xcoldgass2017,xgass2018,piotrowska2022}.

\begin{figure}
\centering
\includegraphics[width=\columnwidth]{../figures/topic-07.pdf}
\caption{SDSS optical denominator/proxy diagnostic for m3\_p2\_gas\_depletion\_efficiency. This is a follow-up target definition or baseline, not a physical-feedback measurement.}
\label{fig:m3-p2-gas-depletion-efficiency}
\end{figure}


\subsection{SDSS target vector for feedback-model validation}
The follow-up goal here is to provide a compact observed target vector for forward modelling, not a direct simulation comparison. What compact SDSS target vector of low-sSFR fraction, optical AGN incidence, and colour versus mass and redshift can be used for forward-model validation? The pilot writes 15 mass-redshift cells with \(n \geq 50\) as a compact validation vector. Across mass bins, low-sSFR fractions span 0.005-0.729, and optical AGN fractions span 0.003-0.520. The output is an observed target vector for simulation forward modelling, not a direct simulation comparison. The missing observables are simulation mocks passed through the SDSS, MaNGA, ALMA, X-ray, and radio selection functions and aperture or noise models \citep{simba2019,tng2019,eagle2015}.

\begin{figure}
\centering
\includegraphics[width=\columnwidth]{../figures/topic-08.pdf}
\caption{SDSS optical denominator/proxy diagnostic for m3\_p3\_simulation\_validation. This is a follow-up target definition or baseline, not a physical-feedback measurement.}
\label{fig:m3-p3-simulation-validation}
\end{figure}


\section{Package decision}
These eight notes should remain supplementary until the missing observables are added. They are suitable as follow-up target definitions, denominator baselines, or appendix material under the main result, but not as independent causal-feedback papers in their current SDSS-only form.

\section{Local reproducibility}
This PDF was generated from the local candidate package \texttt{RP1\_FLAGSHIP\_WITH\_SUPPLEMENT\_20260709T013510Z}. It does not replace any public-linked PDF and does not touch public pages, live roots, product databases, deployment state, git history, billing/OAuth state, cron jobs, or external submission systems.


\begin{thebibliography}{}
\bibitem[Abdurro'uf et al.(2022)]{sdssdr17} Abdurro'uf, Accetta, K., Aerts, C., et al. 2022, ApJS, 259, 35
\bibitem[Baldwin et al.(1981)]{baldwin1981} Baldwin, J.~A., Phillips, M.~M., \& Terlevich, R. 1981, PASP, 93, 5
\bibitem[Best et al.(2005)]{best2005} Best, P.~N., Kauffmann, G., Heckman, T.~M., et al. 2005, MNRAS, 362, 25
\bibitem[Brinchmann et al.(2004)]{brinchmann2004} Brinchmann, J., Charlot, S., White, S.~D.~M., et al. 2004, MNRAS, 351, 1151
\bibitem[Carniani et al.(2017)]{carniani2017} Carniani, S., Marconi, A., Maiolino, R., et al. 2017, A\&A, 605, A42
\bibitem[Catinella et al.(2018)]{xgass2018} Catinella, B., Saintonge, A., Janowiecki, S., et al. 2018, MNRAS, 476, 875
\bibitem[Cicone et al.(2014)]{cicone2014} Cicone, C., Maiolino, R., Sturm, E., et al. 2014, A\&A, 562, A21
\bibitem[Dave et al.(2019)]{simba2019} Dave, R., Angles-Alcazar, D., Narayanan, D., et al. 2019, MNRAS, 486, 2827
\bibitem[Dekel \& Birnboim(2006)]{dekel2006} Dekel, A., \& Birnboim, Y. 2006, MNRAS, 368, 2
\bibitem[Fabian(2012)]{fabian2012} Fabian, A.~C. 2012, ARA\&A, 50, 455
\bibitem[Fiore et al.(2017)]{fiore2017} Fiore, F., Feruglio, C., Shankar, F., et al. 2017, A\&A, 601, A143
\bibitem[Heckman \& Best(2014)]{heckmanbest2014} Heckman, T.~M., \& Best, P.~N. 2014, ARA\&A, 52, 589
\bibitem[Kauffmann et al.(2003a)]{kauffmann2003bpt} Kauffmann, G., Heckman, T.~M., Tremonti, C., et al. 2003a, MNRAS, 346, 1055
\bibitem[Kewley et al.(2001)]{kewley2001} Kewley, L.~J., Dopita, M.~A., Sutherland, R.~S., Heisler, C.~A., \& Trevena, J. 2001, ApJ, 556, 121
\bibitem[Kewley et al.(2006)]{kewley2006} Kewley, L.~J., Groves, B., Kauffmann, G., \& Heckman, T. 2006, MNRAS, 372, 961
\bibitem[LaMassa et al.(2013)]{lamassa2013} LaMassa, S.~M., Heckman, T.~M., Ptak, A., \& Urry, C.~M. 2013, ApJL, 765, L33
\bibitem[McNamara \& Nulsen(2007)]{mcnamara2007} McNamara, B.~R., \& Nulsen, P.~E.~J. 2007, ARA\&A, 45, 117
\bibitem[Nelson et al.(2019)]{tng2019} Nelson, D., Springel, V., Pillepich, A., et al. 2019, Computational Astrophysics and Cosmology, 6, 2
\bibitem[Peng et al.(2010)]{peng2010} Peng, Y.-j., Lilly, S.~J., Kovac, K., et al. 2010, ApJ, 721, 193
\bibitem[Piotrowska et al.(2022)]{piotrowska2022} Piotrowska, J.~M., Bluck, A.~F.~L., Maiolino, R., \& Peng, Y.-j. 2022, MNRAS, 512, 1052
\bibitem[Saintonge et al.(2017)]{xcoldgass2017} Saintonge, A., Catinella, B., Tacconi, L.~J., et al. 2017, ApJS, 233, 22
\bibitem[Schaye et al.(2015)]{eagle2015} Schaye, J., Crain, R.~A., Bower, R.~G., et al. 2015, MNRAS, 446, 521
\bibitem[Stasinska et al.(2008)]{stasinska2008} Stasinska, G., Asari, N.~V., Cid Fernandes, R., et al. 2008, MNRAS, 391, L29
\bibitem[Stasinska et al.(2015)]{stasinska2015} Stasinska, G., Costa Duarte, M.~V., Vale Asari, N., Cid Fernandes, R., \& Sodre, L. 2015, MNRAS, 449, 559
\bibitem[Veilleux et al.(2005)]{veilleux2005} Veilleux, S., Cecil, G., \& Bland-Hawthorn, J. 2005, ARA\&A, 43, 769
\bibitem[Wetzel et al.(2013)]{wetzel2013} Wetzel, A.~R., Tinker, J.~L., Conroy, C., \& van den Bosch, F.~C. 2013, MNRAS, 432, 336
\bibitem[York et al.(2000)]{york2000} York, D.~G., Adelman, J., Anderson, J.~E., Jr., et al. 2000, AJ, 120, 1579
\end{thebibliography}

\end{document}
